% =============================================================================
% APPENDICE: La bussola in 7 domande
% Versione stampabile autonoma: misc/bussola_stampabile.tex
% =============================================================================

\chapter{La bussola in 7 domande}
\label{app:bussola}

Questo libro non si chiude con un elenco di esercizi, ma con sette domande. Non servono a diventare immuni dai bias, cosa che nessuno sa fare. Servono a rallentare nel punto giusto: davanti a una notizia, a una decisione importante, a una risposta troppo convincente comparsa su uno schermo. Non vanno usate tutte, sempre: per le scelte di ogni giorno ne basta una. Per quelle che contano, vale la pena passarle tutte.

\vspace{1em}

\begin{tcolorbox}[enhanced, colback=white, colframe=modernblack, boxrule=1pt, arc=2mm, left=10pt, right=10pt, top=8pt, bottom=8pt]
\begin{center}
{\sffamily\Large\bfseries La bussola in 7 domande}\\[0.3em]
{\sffamily\small\color{moderngray}da \textit{\booktitle} --- \authorname}
\end{center}
\vspace{0.4em}
% Contenuto della «Bussola in 7 domande».
% Usato sia nell'appendice del libro (appendix/bussola.tex)
% sia nella versione stampabile autonoma (misc/bussola_stampabile.tex).
\begin{enumerate}[label={\sffamily\bfseries\large\arabic*}, leftmargin=2.2em, itemsep=0.9em]
	\item \textbf{Prove contrarie.} \textit{Che cosa mi farebbe cambiare idea?} Se non riesco a immaginare nessuna prova capace di smentirmi, non sto ragionando: sto difendendo.

	\item \textbf{Probabilità.} \textit{Quanto è frequente, in generale, quello che sto valutando?} Prima i numeri, poi la storia: il caso che mi colpisce somiglia allo stereotipo, ma quante volte accade davvero?

	\item \textbf{Fonte.} \textit{Da dove viene questa informazione, e come lo saprebbe?} Uno studio, un'esperienza personale, un'opinione, un testo generato da una macchina? Una risposta fluida non è una risposta verificata.

	\item \textbf{Incentivo.} \textit{Chi guadagna se ci credo, o se la condivido?} Vale per chi mi parla, per la piattaforma che me la mostra e per me stesso: a volte l'incentivo più forte è il desiderio di avere ragione.

	\item \textbf{Costo dell'errore.} \textit{Che cosa succede se sbaglio, e quanto è reversibile?} Un errore su un film costa poco; uno sulla salute, sul denaro o su una persona può costare molto. Più alta la posta, più lenta la decisione.

	\item \textbf{Tempo.} \textit{Devo decidere adesso?} Stanchezza, fretta e indignazione sono le condizioni in cui le scorciatoie vincono. Se posso, rimando di un'ora, di una notte.

	\item \textbf{Confronto.} \textit{Chi, tra le persone che rispetto, la vede diversamente?} Chiedo il suo argomento più forte, non la sua versione più facile da smontare. È la correzione che da solo non vedrei.
\end{enumerate}

\end{tcolorbox}

\vspace{1em}

\noindent{\small Questa pagina esiste anche in versione stampabile autonoma (file \texttt{docs/bussola-7-domande.pdf} nel repository del libro, generato da \texttt{misc/bussola\_stampabile.tex}), da staccare e tenere vicino alla scrivania.}
